\documentclass{article}
\usepackage[UTF8]{ctex}
\usepackage{amsmath,amssymb,amsfonts}
\usepackage{enumitem}
\usepackage{geometry}
\geometry{a4paper, margin=1in}
% === 错题本排版优化（不参与题目标识） ===
\setlength{\parindent}{0pt}
\setlength{\parskip}{0.35em plus 0.10em minus 0.08em}
\setlength{\emergencystretch}{3em}
\setlist{itemsep=0.25em, topsep=0.40em, parsep=0pt}
\allowdisplaybreaks[2]
% === 错题本排版优化结束 ===
\begin{document}
\section*{山东大学2024年数学分析真题}
\section*{121．山东大学2024年数学分析试卷}

\subsection{题目1}
1．求 \(\sum_{n=1}^{\infty} \frac{x+n(-1)^{n}}{x^{2}+n^{2}}\) 的条件收敛域，绝对收敛域，以及一致收敛域。

\begin{itemize}

\subsection{题目2}
\end{itemize}

\subsection{题目2}
原题边界“2．”：比较下列无穷大量。

\begin{itemize}
\item[（1）] \(x\) 与 \((\ln x)^{100}(x \rightarrow+\infty)\) ．
  \item[（2）] \((\ln x)^{100}\) 与 \(\mathrm{e}^{(\ln x) 100}(x \rightarrow+\infty)\) ．
\end{itemize}

\subsection{题目3}
3．计算第一类曲面积分 \(\iint_{\sum}(x+y+z) \mathrm{d} S\) ，其中曲面 \(\sum\) 是左半球面 \(x^{2}+y^{2}+z^{2}=1, y \leq 0\) ．

\subsection{题目4}
4．计算 \(I(\alpha)=\int_{0}^{+\infty} \frac{\arctan (\alpha x)}{x\left(1+x^{2}\right)} \mathrm{d} x(\alpha>0)\) ．

\subsection{题目5}
1．设 \(f_{n}(x)=1+\frac{1}{2^{x}}+\frac{1}{3^{x}}+\cdots+\frac{1}{n^{x}}\) ，证明：对任何 \(\alpha>1\) ，函数列 \(\left\{f_{n}(x)\right\}\) 在 \([\alpha,+\infty)\) 上一致收敛。

\subsection{题目6}
2．解答如下问题：

\begin{itemize}
  \item[（1）] 证明不等式 \(\frac{x}{1+x}<\ln (1+x)<x(x>0)\) ．
  \item[（2）] 利用上述不等式证明数列 \(\left\{a_{n}\right\}\) 存在极限，其中 \(a_{n}=1+\frac{1}{2}+\frac{1}{3}+\cdots+\frac{1}{n}-\ln n\) ．
\end{itemize}

\subsection{题目7}
3．设 \(f(x)\) 在区间［0，1］上连续，证明：

\[
\int_{0}^{\pi} x f(\sin x) \mathrm{d} x=\frac{\pi}{2} \int_{0}^{\pi} f(\sin x) \mathrm{d} x .
\]

并由此计算 \(I=\int_{0}^{\pi} x \frac{\sin ^{2 n} x}{\sin ^{2 n} x+\cos ^{2 n} x} \mathrm{~d} x\) ．

\subsection{题目8}
1．设 \(z=z(x, y)\) 满足 \(F\left(x+\frac{z}{y}, y+\frac{z}{x}\right)=0\) ，其中 \(z(x, y), F(u, v)\) 均可微，求 \(z-x z_{x}^{\prime}-y z_{y}^{\prime}\) ．

\subsection{题目9}
2．设二元连续可微函数 \(F\) 在直角坐标下可写为 \(F(x, y)=f(x) g(y)\) ，在极坐标下可写为\(F(r \cos \theta, r \sin \theta) . h(r)\) ，若 \(F(x, y)\) 无零点，求 \(F(x, y)\) ．
\end{document}
